\section{Scheme, normalization, and the complete two-loop atlas}
\label{sec:scheme}

We work with one real scalar field in massive Euclidean $\phi^3$ theory in
six dimensions. Every internal line has the same mass squared $M=m^2>0$,
and
\[
 D(r)=\frac{1}{r^2+M},\qquad
 \int_k=\int_{\mathbb R^6}\frac{d^6k}{\pi^3},\qquad
 \langle r\rangle=\sqrt{r^2+m^2}.
\]
Unless a weight is explicitly attached, $J_R(G)=\int_k\int_q RI_G$
denotes an unweighted scalar graph integral. It excludes the fixed-external-label
symmetry weight $1/S_G$, coupling factors, and conventions for the signs of
effective-action vertices. Replacing each normalized loop measure by
$d^6k/(2\pi)^6$ multiplies a two-loop result by $(4\pi)^{-6}$.
The graph atlas consists of connected amputated 1PI multigraphs, with fixed
external labels and with tadpoles retained.

\subsection{Recursive BPH and the Taylor conditions}

For a graph with $E$ external legs, cubic valence implies the superficial
degree
\[
 \omega(G)=6L-2I=6-2E.
\]
The operator $t_G^{\omega(G)}$ takes the Taylor polynomial, through total
degree $\omega(G)$, in independent physical external momentum components
at zero, holding loop variables fixed. A subgraph Taylor operator similarly
acts on all of that subgraph's independent external momenta, including
momenta which are loop variables of the surrounding graph.

In the Bogoliubov recursion, write
\begin{align*}
 C(\gamma)&=-t_\gamma^{\omega(\gamma)}\overline R I_\gamma,\\
 \overline R I_G&=I_G+
 \sum_{\substack{\mathcal S\ne\varnothing\\
                 \mathcal S\text{ a disjoint proper spinney}}}
 I_{G/\mathcal S}\prod_{\gamma\in\mathcal S} C(\gamma),\\
 RI_G&=(1-t_G^{\omega(G)})\overline R I_G.
\end{align*}
Contractions and local polynomial insertions are understood in the second
line. Counterterms of nested subgraphs are themselves recursively prepared;
overlapping proper subgraphs cannot occur together in one disjoint spinney.
In particular, the two overlapping triangles of the two-point diamond give
two single-subgraph terms, with no product of their counterterms.
For the present atlas the overall Taylor orders are $4,2,0$ for $E=1,2,3$;
the proper bubble and triangle orders are respectively $2$ and $0$.

The terminology is BPH throughout. A primary account of the recursive
framework is A.~D.~Kennedy, \emph{A Simple Proof of the BPH Theorem}
(1996), \texttt{arXiv:hep-th/9612113}, especially its discussion of Henge
and Bogoliubov recursion in Section~3.2. A forest expansion is a useful
finite combinatorial comparison. The graph-specific parameter formulas,
explicit constants, and momentum decompositions below are derived here;
they are not attributed to Kennedy's paper. Nor does this finite atlas
constitute a new proof of an arbitrary-loop BPH theorem.

\subsection{Common regulator and the meaning of an ordinary finite integral}

Use, on every original edge, the same Schwinger lower cutoff
\[
 D_\eta(r)=\int_\eta^\infty
             e^{-\alpha(r^2+M)}\,d\alpha,\qquad \eta>0.
\]
Taylor subtraction is performed on these regulated propagators, before
Gaussian integration. Each counterterm retains the original edge labels and
the common domain $\alpha_e\ge\eta$. Thus every regulated term is integrable;
Fubini and differentiation may be used there. If
$\alpha_e=\rho x_e$, $\sum_e x_e=1$, the common radial lower limit is
\[
 T(x)=\frac{\eta}{\min_e x_e}.
\]
It is independent of external momenta. One must not replace a regulated
combination by separate unregulated integrals of its divergent terms.

For any unit direction, the absolute derivative bounds
\[
 |\partial^nD_\eta(r)|\le c_n\langle r\rangle^{-2-n},\qquad
 (c_0,c_1,c_2,c_3)=(1,2,10,72)
\]
are uniform in $\eta\ge0$. To see this, differentiate the exponential;
each term is bounded by a positive Schwinger moment. The integral from
$\eta$ to infinity is at most the corresponding integral from zero to
infinity. For example, the second derivative is bounded by
$2\langle r\rangle^{-4}+8r^2\langle r\rangle^{-6}$, and the third by
$24|r|\langle r\rangle^{-6}+48|r|^3\langle r\rangle^{-8}$.
The shift comparison
\[
 \langle r+a\rangle\ge
 \frac{\langle r\rangle}{1+|a|/m}
\]
follows from the triangle inequality applied to $(r,m)$.
These facts also bound propagator differences by the mean value formula.
They supply regulator-independent momentum majorants in the chapters below.

There are two distinct steps in each equivalence proof. A momentum-space
$L^1$ majorant proves $\int RI_\eta\to\int RI$ for the original ordinary
integral. A separate parameter-space majorant proves that the same regulated
left side converges to the displayed finite parameter formula. Identifying
these limits is legitimate; inferring original momentum $L^1$ solely from a
post-Gaussian parameter norm would not be.

\subsection{One-point graph and the finite linear condition}

For the sole two-loop one-point graph, momentum conservation fixes its only
physical external momentum to zero. There is no independent external
variable left in the prepared integrand. Consequently the overall Taylor
operator is the identity, even though its nominal order is four:
\[
 t_G^4\overline R I_{G,\eta}=\overline R I_{G,\eta},\qquad
 RI_{G,\eta}=0,
 \qquad J_R(\mathrm{L2E1G1})=0.
\]
This holds before any divergent term is integrated. The prepared graph may
contain nontrivial proper-subgraph counterterms, but the identity applies
to their complete sum. In this full zero-momentum Taylor scheme the local
linear counterterm has coefficient
$-\overline R J_{G,\eta}(0)$, with the sign translated according to the
chosen action convention. Its finite part fixes the graphwise renormalized
one-point amplitude to zero. The bare tadpole is not asserted to vanish.
This is a renormalization condition on the linear term; it is not an
independent proof of a physical vacuum expectation value condition.

\subsection{All ten fixed-label graphs and their scalar kernels}

Let $s=p^2$ in a two-point graph and let
$s_i=p_i^2$, $p_1+p_2+p_3=0$, in a three-point graph. Denote the diamond
kernel by $J_{\mathrm{ov}}(s)$, the primitive kernel by
$J_{\mathrm{pr}}(s_1,s_2,s_3)$, and the standard triangle-insertion kernel
by $J_{\mathrm{tr}}(s_1,s_2,s_3)$, where its first argument is the leg
attached to the subdivergent triangle. For a bubble insertion write
$J_{\mathrm{b},2}(s)$ or $J_{\mathrm{b},3}(a,b,c)$; in the latter,
$a,b$ are adjacent to the insertion and $c$ is opposite it.
All entries below use the same $M$.

\begin{center}
\small
\begin{tabular}{llll}
Graph & scalar kernel & proper subtraction & $S_G$\\ \hline
L2E1G1 & $0$ & bubble and triangles & $4$\\
L2E2G1 & $J_{\mathrm{ov}}(s)$ & two overlapping triangles & $2$\\
L2E2G2 & $J_{\mathrm{b},2}(s)$ & one bubble & $2$\\
L2E3G1 & $J_{\mathrm{pr}}(s_1,s_2,s_3)$ & none & $2$\\
L2E3G2 & $J_{\mathrm{tr}}(s_1,s_2,s_3)$ & triangle on leg $1$ & $1$\\
L2E3G3 & $J_{\mathrm{tr}}(s_2,s_1,s_3)$ & triangle on leg $2$ & $1$\\
L2E3G4 & $J_{\mathrm{b},3}(s_1,s_2,s_3)$ & bubble opposite leg $3$ & $2$\\
L2E3G5 & $J_{\mathrm{tr}}(s_3,s_1,s_2)$ & triangle on leg $3$ & $1$\\
L2E3G6 & $J_{\mathrm{b},3}(s_1,s_3,s_2)$ & bubble opposite leg $2$ & $2$\\
L2E3G7 & $J_{\mathrm{b},3}(s_2,s_3,s_1)$ & bubble opposite leg $1$ & $2$
\end{tabular}
\end{center}
These mappings use the actual fixed external labels. A symmetric momentum
point cannot test them: adjacent-leg exchange is a symmetry of the standard
insertion kernels, whereas exchange with their distinguished leg generally
changes the answer. The primitive kernel is symmetric under all three
external-leg permutations.

\subsection{Physical input domain and the scope of the bounds}

The mathematical three-point domain is
\[
 s_i\ge0,\qquad
 G=4s_1s_2-(s_3-s_1-s_2)^2\ge0.
\]
It is invariant under all leg permutations, and includes degenerate and
zero momenta. In the software, converted binary64 values are checked using
exact rational arithmetic. A negative requested $G$ is tolerated only when
its defect is at most $10^{-14}$ times the smallest of the three Cauchy
product scales $4ab+(c-a-b)^2$. The largest invariant is then moved inward
by binary64 steps until the effective triple has exact $G\ge0$. Requested
and effective inputs are recorded separately; all kernels and bounds use
the effective physical triple. In particular a zero leg requires equality
of the other two inputs, and an absolute error tolerance cannot hide an
inconsistent tiny leg. This numerical boundary convention does not weaken
the exact mathematical hypothesis above.

All compact-uniform statements mean that physical external momenta have a
fixed upper bound and $m\ge m_0>0$. The stronger global logarithmic scale
bounds are stated separately. Analytic amplitude bounds, original momentum
$L^1$ bounds, and quadrature refinement differences are different objects;
none is silently substituted for another. The scope is these ten massive
Euclidean graphs in the declared finite scheme, not a massless limit,
Minkowski continuation, general finite-scheme conversion, or an arbitrary-loop
induction theorem.
